% Generated by tools/materialize_round8_referee.py; do not hand-edit.
\section{Round-eight reconstruction: edge-flow compactification and projective epigraph recovery}
\label{sec:r8-a3}

The state now retains the inducing branch, its transition edge, and the actual
excursion profile.  Recession is obtained from a directed family of genuine
finite-coordinate LDPs; no full-sequence Gamma limit is declared in advance.

\subsection{The exact marked edge-flow state}

Let $\mathcal G=(\mathcal V,\mathcal E)$ be the countable first-return graph.
An admissible path is a sequence of edges $e_j=(a_j,a_{j+1})$.  For an actual
excursion point $y_j\in Y_{a_j}$ put
\[
 K(y_j,a_j)=\frac1{r(a_j)}\sum_{\ell<r(a_j)}
 \delta_{\Theta^\ell\Gamma(y_j)}.
\]
For $R_N=\sum_{j<N}r(a_j)$ define
\[
 M_N=R_N^{-1}\sum_{j<N}r(a_j)\delta_{(a_j,a_{j+1})},
\qquad
 \Xi_N=R_N^{-1}\sum_{j<N}r(a_j)
 \delta_{(a_j,a_{j+1},K(y_j,a_j))}.
\]
The mark contains the actual profile, not its branch conditional expectation.

Choose finite edge sets $\mathcal E_m\uparrow\mathcal E$ and finitely many
continuous profile tests $\phi_1,\ldots,\phi_m$.  The projection $\pi_m\Xi$
records the masses of edges in $\mathcal E_m$, their incoming and outgoing
marginals, the first $m$ profile moments, and the escaped clock
\[
 z_m=1-\sum_{e\in\mathcal E_m}\Xi(e,\mathcal Z).
\]

\begin{lemma}[Rate data are functions of the state]
\label{lem:r8-a3-state}
The edge marginals of $\Xi$ determine admissibility, stationarity, the
length-weighted entropy flow, and the induced potential average.  Distinct
branches with the same return length and profile remain distinct coordinates.
\end{lemma}

\begin{proof}
Incoming and outgoing marginals are read from the first two components of the
mark.  Stationarity is equality of these marginals.  Conditional transition
probabilities are the Radon--Nikodym ratios of the edge flow with respect to its
incoming marginal, so their entropy is determined by $\Xi$.  The branch
potential and length are edge labels and are integrated directly.  No
injectivity of $a\mapsto(r(a)^{-1},K_a)$ is used.
\end{proof}

\subsection{Finite projections and the projective epigraph theorem}

Let $I_m$ be the good finite-dimensional rate obtained by the exponential
change of measure on $\pi_m\Xi_N$.  The source class contains all edge
indicators in $\mathcal E_m$ and all profile tests $\phi_1,\ldots,\phi_m$.
Consistency gives
\[
 I_m(z)=\inf\{I_{m+1}(z'): \pi_m z'=z\}.
\]

\begin{theorem}[Projective epigraph LDP]
\label{thm:r8-a3-projective}
Assume exponential tightness of the return-length moment and the finite-source
coefficient theorem of A2.  Then $\Xi_N$ satisfies a full good LDP on the
projective state space with rate
\[
 I(\Xi)=\sup_m I_m(\pi_m\Xi).
\]
Every finite-rate state has a diagonal sequence of actual admissible paths
whose costs converge to $I(\Xi)$.
\end{theorem}

\begin{proof}
The upper bound follows from exponential tightness and the finite-projection
upper bounds.  For the lower bound, choose $m_j\uparrow\infty$ so that
$I_{m_j}(\pi_{m_j}\Xi)\uparrow I(\Xi)$.  The finite-dimensional lower theorem
provides an actual admissible path block in a $2^{-j}$ neighborhood of
$\pi_{m_j}\Xi$ with cost at most $I_{m_j}+2^{-j}$.  Concatenate these blocks as successive excursion sequences.  If the terminal
edge of one block does not meet the initial edge of the next, insert a finite
connecting excursion path supplied by transitivity of the return graph.
Every intermediate base return is a legal edge-flow transition rather than a
forbidden connector inside one first-return word.  Choose block multiplicities
increasing so fast that all connecting paths and earlier blocks have vanishing
clock fraction.  The diagonal path converges in every projection and has the desired
cost.  This proves the recovery statement rather than assuming it.
\end{proof}

\subsection{Recession as boundary edge flow}

The boundary mass $z_m$ is monotone in $m$ and converges to
$z_\infty\in[0,1]$.  A boundary profile is a projective limit of the moments
carried by edges escaping every finite set.  Define
\[
 I_{\rm rec}(z)=\sup_m I_m^{\rm boundary}(\pi_mz).
\]

\begin{theorem}[Actual-excursion recession recovery]
\label{thm:r8-a3-recession}
The recession functional is lower semicontinuous, convex under concatenation of
excursion blocks, and has the dual formula
\[
 I_{\rm rec}(z)=\sup_{F\in C_{\rm cyl}}
 \{z(F)-q_\infty(F)+q_\infty(0)\}.
\]
For every finite-rate $z$ there are actual long first-return excursions whose
finite profile projections approach $z$ and whose logarithmic costs approach
$I_{\rm rec}(z)$.
\end{theorem}

\begin{proof}
Each finite boundary rate is a closed convex conjugate and the directed
supremum is lower semicontinuous.  The recovery part of
Theorem~\ref{thm:r8-a3-projective}, applied to states with clock mass outside
$\mathcal E_m$, produces actual excursions at each finite stage.  A diagonal
selection makes the return length tend to infinity while matching every
profile projection.  Convex combinations are realized by consecutive long
excursions separated by genuine base returns; no claim is made that their
concatenation is one first-return word.  The cylinder dual determines the
projective supremum by separation.
\end{proof}

\subsection{Moment closure without an unbounded reciprocal coordinate}

For every empirical state,
\[
 \Xi_N(1)=1.
\]
The recurrent clock in a finite set and its complement add exactly to one.
The boundary coordinate is therefore a continuous ledger coordinate, not the
weak limit of the unbounded function $x^{-1}$.

\begin{lemma}[Exact recurrent--defect normalization]
\label{lem:r8-a3-normalization}
Every projective limit satisfies
\[
 m_{\rm rec}(1)+z_\infty=1.
\]
Conversely every stationary edge flow with this identity is approximated by
normalized empirical edge flows.
\end{lemma}

\begin{proof}
For each $m$ the identity
$m_{N}(\mathcal E_m)+z_{N,m}=1$ is exact and involves bounded coordinates.
Pass first to $N\to\infty$ and then $m\to\infty$.  For the converse, truncate
the flow, decompose the finite circulation into cycles, and append boundary
excursion blocks supplied by Theorem~\ref{thm:r8-a3-recession}.  Normalize by
total clock.
\end{proof}

\subsection{Finite-core approximation and entropy}

\begin{theorem}[Stationary flow approximation]
\label{thm:r8-a3-flow-approx}
Every finite-rate stationary countable edge flow is approximated by stationary
flows on finite strongly connected subgraphs, with convergence of potential,
length moment, and entropy.  Periodic components retain their period.
\end{theorem}

\begin{proof}
Truncate to edges carrying all but $\varepsilon$ of the length, potential, and
information moments.  The imbalance at the retained vertices has total mass
$O(\varepsilon)$.  Decompose the original stationary flow into directed cycles
and keep finitely many cycles that repair this imbalance.  Their union is a
finite graph; apply the construction separately to each cyclic class instead
of adding a period-breaking orbit.  The entropy error is controlled by the
truncated information moment and tends to zero.
\end{proof}

\subsection{Spectral versus recession sectors}

\begin{theorem}[SPR/recession dichotomy]
\label{thm:r8-a3-dichotomy}
A strongly positive recurrent component contributes through its isolated
renewal pole and the recurrent edge-flow rate.  If the pole is lost at the
critical pressure, its normalized mass escapes every finite edge set and is
represented by $I_{\rm rec}$.  No component is counted in both sectors.
\end{theorem}

\begin{proof}
On a finite core the renewal resolvent is compact and has an isolated Perron
pole.  Uniform return-length integrability keeps this pole isolated in the
countable limit, which is the SPR case.  If uniform integrability fails, any
normalized approximate eigenvector has vanishing mass on every finite edge set;
its finite projections converge to a boundary state.  The projective epigraph
rate of that state is exactly the recession cost.  The alternatives are
mutually exclusive by tightness.
\end{proof}

\subsection{Exact physical time}

For a deterministic time $T$, keep the terminal cut excursion
$(a_*,u_*,K_*^{\rm cut})$, where $u_*\in[0,1]$ is the used fraction of its roof.
This coordinate is part of the state rather than discarded.

\begin{theorem}[Deterministic-time path LDP]
\label{thm:r8-a3-time}
The edge-flow LDP contracts to the full empirical path at every deterministic
physical time.  The terminal cut coordinate gives exact time $T$ and its cost
is the corresponding fractional excursion cost.  When the terminal roof is
$o(T)$ it is exponentially equivalent to zero; when it is of order $T$ it is
already represented by the recession sector.
\end{theorem}

\begin{proof}
Stop at the first return whose accumulated roof exceeds $T$ and retain the
used fraction of that excursion.  This makes the clock identity exact.
Exponential tightness separates the event $r(a_*)\le\delta T$, whose profile
error is at most $\delta$, from $r(a_*)>\delta T$, which is a boundary
excursion and is governed by Theorem~\ref{thm:r8-a3-recession}.  Let
$\delta\downarrow0$ after the LDP bounds.
\end{proof}
